% ---
% Reemplaza líneas 1907-1967 de Docs/Technical_Report.tex
% Decisiones: D-018, D-036
% ---

\subsubsection{Métricas de clasificación}

El desempeño del clasificador MLP sobre cada condición experimental se evalúa a partir de
las predicciones de probabilidad de malignidad $\hat{y} \in [0, 1]$ generadas por la función
sigmoide de salida. Para las métricas dependientes de umbral, se adopta el umbral de decisión
operativo estándar de 0.5 ($\hat{y} \geq 0.5 \implies \text{maligno}$):

\textbf{Exactitud (\textit{Accuracy}):}
\begin{equation}
  \text{Acc} =
  \frac{TP + TN}{TP + TN + FP + FN}
  \label{eq:accuracy}
\end{equation}

\textbf{Precisión (\textit{Precision}):}
\begin{equation}
  \text{Prec} =
  \frac{TP}{TP + FP}
  \label{eq:precision}
\end{equation}

\textbf{Exhaustividad (\textit{Recall}):}
\begin{equation}
  \text{Rec} =
  \frac{TP}{TP + FN}
  \label{eq:recall}
\end{equation}

\textbf{F1-Score:}
\begin{equation}
  F_1 =
  \frac{2 \cdot \text{Prec} \cdot \text{Rec}}
  {\text{Prec} + \text{Rec}}
  \label{eq:f1}
\end{equation}

\textbf{Área bajo la curva ROC (AUC-ROC):}
\begin{equation}
  \text{AUC} =
  \int_0^1 \text{TPR}(t)\,d\,\text{FPR}(t)
  \label{eq:auc_roc}
\end{equation}

\noindent donde $TP$, $TN$, $FP$ y $FN$ son los
verdaderos positivos, verdaderos negativos, falsos
positivos y falsos negativos respectivamente;
$\text{TPR}(t) = TP/(TP+FN)$ es la tasa de verdaderos
positivos y $\text{FPR}(t) = FP/(FP+TN)$ es la tasa
de falsos positivos, ambas evaluadas al umbral de
decisión $t$.

El \textbf{Área bajo la curva ROC (AUC-ROC)}, formalizada en la ecuación~\eqref{eq:auc_roc},
constituye la métrica principal de comparación técnica entre condiciones. Al integrar la
tasa de verdaderos positivos en función de la tasa de falsos positivos sobre la totalidad de
los umbrales posibles $t \in [0, 1]$, el AUC-ROC cuantifica la capacidad discriminativa intrínseca
del modelo con independencia del umbral operativo y de forma robusta frente al desbalance
de clases del CBIS-DDSM.

El protocolo de validación evalúa todas las métricas mediante validación cruzada de 5 \textit{folds}
estratificada por clase y agrupada por paciente sobre el conjunto de entrenamiento, así
como mediante una evaluación final sobre el conjunto de prueba oficial, que se usa una sola vez por modelo, reportando
intervalos de confianza al 95\,\% derivados de 2\,000 remuestreos \textit{bootstrap} pareados por
paciente.
